\section{Distribusi binomial}
\label{binomialModel}

\index{distribution!binomial|(}

\termsub{Distribusi binomial}{distribution!binomial}
digunakan untuk menggambarkan banyaknya sukses ketika jumlah percobaannya
tetap.
%,
%and this distribution is occasionally used in statistics,
%especially when doing more careful analysis of samples
%of data where simpler tools are not helpful.
Distribusi ini berbeda dari distribusi geometrik, yang menggambarkan
banyaknya percobaan yang harus kita tunggu sebelum mengamati suatu sukses.


\subsection{Distribusi binomial}

%\newcommand{\insureSprob}{0.7}
%\newcommand{\insureSperc}{70\%}
%\newcommand{\insureFprob}{0.3}
%\newcommand{\insureFperc}{30\%}
%\newcommand{\insureDistA}{0.7}
%\newcommand{\insureDistB}{0.21}
%\newcommand{\insureDistC}{0.063}
%\newcommand{\insureDistD}{0.019}
%\newcommand{\insureDistE}{0.006}
%\newcommand{\insureCDistA}{0.7}
%\newcommand{\insureCDistB}{0.91}
%\newcommand{\insureCDistC}{0.973}
%\newcommand{\insureCDistCComplement}{0.027}
%\newcommand{\insureCDistD}{0.992}
%\newcommand{\insureCDistE}{0.998}
%\newcommand{\insureGeomMean}{1.43}
\newcommand{\insureS}{\resp{tidak}}
\newcommand{\insureF}{\resp{melampaui}}
% Doesn't consider binomial coefficient in next calculated value.
\newcommand{\insureBinomCinDSingleScenario}{0.103}
\newcommand{\insureBinomCinD}{0.412}
\newcommand{\insureBinomEinHSingleScenario}{0.00454}
\newcommand{\insureBinomEinH}{0.254}
\newcommand{\insureBinomFourtyExpValue}{28}
\newcommand{\insureBinomFourtySD}{2.9}
\newcommand{\insureBinomFourtyLower}{22}
\newcommand{\insureBinomFourtyUpper}{34}

\noindent%
Mari kita bayangkan lagi bahwa kita berada di perusahaan asuransi,
tempat \insureSperc{} tertanggung tidak melampaui batas risiko sendiri.

\begin{examplewrap}
\begin{nexample}{Misalkan perusahaan asuransi mempertimbangkan
    sampel acak berisi empat orang yang diasuransikannya.
    Berapa peluang bahwa tepat satu orang melampaui batas risiko sendiri,
    sedangkan tiga orang lainnya tidak? Agar mudah, kita namai keempat
    orang tersebut Ariana ($A$), Brittany ($B$), Carlton ($C$), dan
    Damian ($D$).}
  \label{insureOneOfFourExceedsDeductible}%
  Mari kita tinjau satu skenario ketika satu orang melampaui batas risiko
  sendiri:
  \begin{align*}
  &P(A=\text{\insureF{}},
      \text{ }B=\text{\insureS{}},
      \text{ }C=\text{\insureS{}},
      \text{ }D=\text{\insureS{}}) \\
    &\quad = P(A=\text{\insureF{}})\ 
        P(B=\text{\insureS{}})\ 
        P(C=\text{\insureS{}})\ 
        P(D=\text{\insureS{}}) \\
    &\quad =  (\insureFprob{})
        (\insureSprob{})
        (\insureSprob{})
        (\insureSprob{}) \\
    &\quad = (\insureSprob{})^3 (\insureFprob{})^1 \\
    &\quad = \insureBinomCinDSingleScenario{}
  \end{align*}
  Namun, terdapat tiga skenario lain: Brittany, Carlton, atau Damian
  dapat menjadi satu-satunya orang yang melampaui batas risiko sendiri.
  Pada masing-masing kasus tersebut, peluangnya kembali sebesar
  $(\insureSprob{})^3 (\insureFprob{})^1$.
  Keempat skenario ini mencakup semua kemungkinan cara tepat satu dari
  keempat orang tersebut melampaui batas risiko sendiri. Karena itu,
  peluang totalnya adalah
  $4 \times (\insureSprob{})^3 (\insureFprob{})^1
      = \insureBinomCinD{}$.
\end{nexample}
\end{examplewrap}

\begin{exercisewrap}
\begin{nexercise}
Verifikasikan bahwa skenario ketika Brittany menjadi satu-satunya orang
yang melampaui batas risiko sendiri memiliki peluang
$(\insureSprob{})^3 (\insureFprob{})^1$.~\footnotemark{}
\end{nexercise}
\end{exercisewrap}
\footnotetext{
  $P(A=\text{\insureS{}},
      \text{ }B=\text{\insureF{}},
      \text{ }C=\text{\insureS{}},
      \text{ }D=\text{\insureS{}})
    = (\insureSprob{})(\insureFprob{})
        (\insureSprob{})(\insureSprob{})
    = (\insureSprob{})^3 (\insureFprob{})^1$.}


Skenario yang diuraikan pada
Contoh~\ref{insureOneOfFourExceedsDeductible} merupakan contoh
eksperimen binomial. \termsub{Distribusi binomial}{distribution!binomial}
menggambarkan peluang untuk memperoleh tepat $k$ sukses dalam $n$
percobaan Bernoulli yang saling independen dengan peluang sukses $p$
(pada Contoh~\ref{insureOneOfFourExceedsDeductible},
$n=4$, $k=3$, dan $p=\insureSprob{}$).
Kita hendak menentukan peluang-peluang yang berkaitan dengan distribusi
binomial secara lebih umum. Dengan kata lain, kita menginginkan rumus
yang menggunakan $n$, $k$, dan $p$ untuk memperoleh peluang tersebut.
Untuk itu, kita menelaah kembali setiap bagian dari
Contoh~\ref{insureOneOfFourExceedsDeductible}.

Terdapat empat orang yang masing-masing dapat menjadi satu-satunya orang
yang melampaui batas risiko sendiri, dan keempat skenario tersebut
memiliki peluang yang sama. Dengan demikian, peluang akhirnya dapat kita
nyatakan sebagai
\begin{align*}
[\text{\# skenario}] \times P(\text{satu skenario})
\end{align*}
Komponen pertama persamaan ini adalah banyaknya cara untuk menempatkan
$k=3$ sukses di antara $n=4$ percobaan. Komponen kedua adalah peluang
salah satu dari keempat skenario yang berpeluang sama tersebut.

\D{\newpage}

Tinjau $P($satu skenario$)$ dalam kasus umum ketika terdapat $k$ sukses
dan $n-k$ kegagalan dalam $n$ percobaan. Pada setiap skenario seperti itu,
kita menerapkan Aturan Perkalian untuk peristiwa yang saling independen:
\begin{align*}
p^k (1 - p)^{n - k}
\end{align*}
Inilah rumus umum kita untuk $P($satu skenario$)$.

Selanjutnya, kita memperkenalkan rumus umum untuk banyaknya cara memilih
$k$ sukses dalam $n$ percobaan, yaitu menempatkan $k$ sukses dan
$n-k$ kegagalan:
\begin{align*}
{n\choose k} = \frac{n!}{k! (n - k)!}
\end{align*}
Besaran ${n\choose k}$ dibaca
\term{n pilih k}.\footnote{Notasi lain untuk $n$ pilih $k$
  mencakup $_nC_k$, $C_n^k$, dan $C(n,k)$.}
Notasi tanda seru (misalnya $k!$) menyatakan ekspresi
\term{faktorial}.\label{factorial_defined}
\begin{align*}
& 0! = 1 \\
& 1! = 1 \\
& 2! = 2\times1 = 2 \\
& 3! = 3\times2\times1 = 6 \\
& 4! = 4\times3\times2\times1 = 24 \\
& \vdots \\
& n! = n\times(n-1)\times...\times3\times2\times1
\end{align*}
Dengan rumus tersebut, kita dapat menghitung banyaknya cara memilih
$k=3$ sukses dalam $n=4$ percobaan:
\begin{align*}
{4 \choose 3} = \frac{4!}{3!(4-3)!}
  = \frac{4!}{3!1!} 
  = \frac{4\times3\times2\times1}{(3\times2\times1) (1)}
  = 4
\end{align*}
Hasil ini tepat sama dengan yang kita peroleh ketika mempertimbangkan
setiap skenario yang mungkin secara saksama pada
Contoh~\ref{insureOneOfFourExceedsDeductible}.

Mengganti banyaknya skenario dengan $n$ pilih $k$ dan peluang satu
skenario dengan $p^k(1-p)^{n-k}$ menghasilkan rumus umum binomial.

\begin{onebox}{Distribusi binomial}
  Misalkan peluang bahwa satu percobaan menghasilkan sukses adalah $p$.
  Maka peluang untuk mengamati tepat $k$ sukses dalam $n$ percobaan yang
  saling independen diberikan oleh\vspace{-1mm}
  \begin{align*}
  {n\choose k}p^k(1-p)^{n-k} = \frac{n!}{k!(n-k)!}p^k(1-p)^{n-k}
  \end{align*}
  Rata-rata, varians, dan simpangan baku banyaknya sukses yang diamati
  adalah\vspace{-2mm}
  \begin{align*}
  \mu &= np
  &\sigma^2 &= np(1-p)
  &\sigma&= \sqrt{np(1-p)}
  \end{align*}
\end{onebox}

\begin{onebox}{Apakah modelnya binomial? Empat syarat yang perlu diperiksa.}
  \label{isItBinomialTipBox}%
  (1) Percobaan-percobaannya saling independen. \\
  (2) Jumlah percobaan, $n$, tetap. \\
  (3) Hasil setiap percobaan dapat diklasifikasikan sebagai
      \emph{sukses} atau \emph{gagal}. \\
  (4) Peluang sukses, $p$, sama untuk setiap percobaan.
\end{onebox}

\D{\newpage}

\begin{examplewrap}
\begin{nexample}{Berapa peluang bahwa 3 dari 8 orang yang dipilih secara
    acak akan melampaui batas risiko sendiri, atau dengan kata lain bahwa
    5 dari 8 orang tidak akan melampauinya? Ingat bahwa
    70\% orang tidak melampaui batas risiko sendiri.}
  Kita hendak menerapkan model binomial, sehingga kita memeriksa
  syarat-syaratnya. Jumlah percobaan tetap ($n=8$) (syarat 2), dan
  setiap hasil percobaan dapat diklasifikasikan sebagai sukses atau gagal
  (syarat 3). Karena sampelnya acak, percobaan-percobaan tersebut saling
  independen (syarat~1), dan peluang sukses sama untuk setiap percobaan
  (syarat~4).

  Pada hasil yang menjadi perhatian, terdapat $k=5$ sukses dalam $n=8$
  percobaan (ingat bahwa sukses adalah seseorang yang \emph{tidak}
  melampaui batas risiko sendiri), dan peluang suksesnya adalah
  $p=\insureSprob{}$. Jadi, peluang bahwa 5 dari 8 orang tidak melampaui
  batas risiko sendiri dan 3 orang melampauinya diberikan oleh
  \begin{align*}
  { 8 \choose 5}(\insureSprob{})^5
  (1-\insureSprob{})^{8-5}
    &= \frac{8!}{5!(8-5)!}
        (\insureSprob{})^5(1-\insureSprob{})^{8-5} \\
    &= \frac{8!}{5!3!}
        (\insureSprob{})^5(\insureFprob{})^3
  \end{align*}
  Untuk menangani bagian faktorial:
  \begin{align*}
  \frac{8!}{5!3!}
    = \frac{8\times7\times6\times5\times4\times3\times2\times1}
        {(5\times4\times3\times2\times1)(3\times2\times1)}
    = \frac{8\times7\times6}{3\times2\times1}
    = 56
  \end{align*}
  Dengan menggunakan $(\insureSprob{})^5(\insureFprob{})^3
    \approx \insureBinomEinHSingleScenario{}$,
  peluang akhirnya sekitar
  $56 \times \insureBinomEinHSingleScenario{}
    \approx \insureBinomEinH{}$.
\end{nexample}
\end{examplewrap}

\begin{onebox}{Menghitung peluang binomial}
  Langkah pertama ketika menggunakan model binomial adalah memeriksa
  bahwa model tersebut tepat. Langkah kedua adalah mengidentifikasi $n$,
  $p$, dan $k$. Pada tahap terakhir, gunakan perangkat lunak atau rumus
  untuk menentukan peluang, lalu tafsirkan hasilnya.%
  \vspace{3mm}

  Jika Anda harus melakukan perhitungan dengan tangan, sering kali
  berguna untuk mencoret sebanyak mungkin faktor yang sama pada pembilang
  dan penyebut koefisien binomial.
\end{onebox}

\begin{exercisewrap}
\begin{nexercise}
Jika kita mengambil sampel acak berisi 40 berkas kasus dari perusahaan
asuransi yang dibahas sebelumnya, berapa banyak kasus yang diperkirakan
tidak melampaui batas risiko sendiri dalam suatu tahun? Berapa simpangan
baku banyaknya kasus yang tidak melampaui batas risiko sendiri?\footnotemark{}
\end{nexercise}
\end{exercisewrap}
\footnotetext{Kita diminta menentukan banyaknya kasus yang diharapkan
  (rata-rata) dan simpangan bakunya. Keduanya dapat dihitung langsung
  dari rumus:
  $\mu = np = 40 \times \insureSprob{}
    = \insureBinomFourtyExpValue$
  dan $\sigma = \sqrt{np(1-p)}
    = \sqrt{40\times \insureSprob{}\times \insureFprob{}}
    = \insureBinomFourtySD{}$.
  Karena secara sangat kasar 95\% pengamatan berada dalam jarak
  2~simpangan baku dari rata-rata
  (lihat Bagian~\ref{variability}), dalam sampel ini mungkin akan kita
  amati setidaknya \insureBinomFourtyLower{} tetapi kurang dari
  \insureBinomFourtyUpper{} orang yang tidak melampaui batas risiko sendiri.}

\begin{exercisewrap}
\begin{nexercise}
Peluang bahwa seorang perokok yang dipilih secara acak akan mengalami
gangguan paru-paru berat sepanjang hidupnya adalah sekitar $0.3$.
Jika Anda memiliki 4 teman yang merokok, apakah syarat-syarat model
binomial terpenuhi?\footnotemark{}
\end{nexercise}
\end{exercisewrap}
\footnotetext{Salah satu jawaban yang mungkin: jika teman-teman tersebut
  saling mengenal, asumsi independensi mungkin tidak terpenuhi. Sebagai
  contoh, orang-orang yang saling mengenal mungkin memiliki kebiasaan
  merokok yang serupa, atau mereka mungkin bersepakat untuk berhenti
  merokok bersama-sama.}

\D{\newpage}

\begin{exercisewrap}
\begin{nexercise}
\label{noMoreThanOneFriendWSevereLungCondition}%
Misalkan keempat teman tersebut tidak saling mengenal dan dapat kita
perlakukan seolah-olah mereka merupakan sampel acak dari populasi.
Apakah model binomial tepat? Berapa peluang bahwa\footnotemark{}
\begin{enumerate}[(a)]
\setlength{\itemsep}{0mm}
\item
    Tidak seorang pun mengalami gangguan paru-paru berat?
\item
    Satu orang mengalami gangguan paru-paru berat?
\item
    Tidak lebih dari satu orang mengalami gangguan paru-paru berat?
\end{enumerate}
\end{nexercise}
\end{exercisewrap}
\footnotetext{Untuk memeriksa apakah model binomial tepat, kita harus
  memverifikasi syarat-syaratnya.
  (i)~Karena kita mengandaikan bahwa teman-teman tersebut dapat
  diperlakukan sebagai sampel acak, mereka saling independen.
  (ii)~Kita memiliki jumlah percobaan yang tetap ($n=4$).
  (iii)~Setiap hasil adalah sukses atau gagal.
  (iv)~Peluang sukses sama pada setiap percobaan karena orang-orang
  tersebut dapat diperlakukan sebagai sampel acak ($p=0.3$ jika kita
  menyebut seseorang yang mengalami gangguan paru-paru sebagai ``sukses'', sebuah pilihan
  istilah yang suram).
  Hitung bagian~(a) dan~(b) dengan rumus binomial:
  $P(0)
    = {4 \choose 0} (0.3)^0 (0.7)^4
    = 1\times1\times0.7^4
    = 0.2401$,
  $P(1)
    = {4 \choose 1} (0.3)^1(0.7)^{3}
    = 0.4116$.
  Catatan: $0!=1$.
  Bagian~(c) dapat dihitung sebagai jumlah bagian~(a) dan~(b):
  $P(0) + P(1) = 0.2401 + 0.4116 = 0.6517$.
  Artinya, terdapat peluang sekitar 65\% bahwa tidak lebih dari satu di
  antara keempat teman Anda yang merokok akan mengalami gangguan
  paru-paru berat.}

\begin{exercisewrap}
\begin{nexercise}
Berapa peluang bahwa sekurang-kurangnya 2 dari 4 teman Anda yang merokok
akan mengalami gangguan paru-paru berat sepanjang hidup mereka?\footnotemark{}
\end{nexercise}
\end{exercisewrap}
\footnotetext{Peluang komplemennya (tidak lebih dari satu orang akan
  mengalami gangguan paru-paru berat) telah dihitung sebagai 0.6517 pada
  Latihan Terbimbing~\ref{noMoreThanOneFriendWSevereLungCondition}.
  Jadi, kita menghitung satu dikurangi nilai tersebut:~0.3483.}

\begin{exercisewrap}
\begin{nexercise}
Misalkan Anda memiliki 7 teman yang merokok dan mereka dapat diperlakukan
sebagai sampel acak perokok.\footnotemark{}
\begin{enumerate}[(a)]
\setlength{\itemsep}{0mm}
\item
    Berapa banyak di antara mereka yang diperkirakan akan mengalami
    gangguan paru-paru berat, yaitu berapa rata-ratanya?
\item
    Berapa peluang bahwa paling banyak 2 dari 7 teman Anda akan mengalami
    gangguan paru-paru berat?
\end{enumerate}
\end{nexercise}
\end{exercisewrap}
\footnotetext{(a)~$\mu=0.3\times7 = 2.1$.
  (b)~$P($0, 1, atau 2 orang mengalami gangguan paru-paru berat$)
      = P(k=0) + P(k=1)+P(k=2) = 0.6471$.}

Selanjutnya, kita meninjau suku pertama dalam peluang binomial, yaitu
$n$ pilih $k$, untuk beberapa skenario khusus.

\begin{exercisewrap}
\begin{nexercise}
Mengapa ${n \choose 0}=1$ dan ${n \choose n}=1$ berlaku untuk sebarang
bilangan $n$?\footnotemark{}
\end{nexercise}
\end{exercisewrap}
\footnotetext{Nyatakan ekspresi-ekspresi ini dengan kata-kata.
  Ada berapa cara berbeda untuk menempatkan 0 sukses dan $n$ kegagalan
  dalam $n$ percobaan? (1 cara.)
  Ada berapa cara berbeda untuk menempatkan $n$ sukses dan 0 kegagalan
  dalam $n$ percobaan? (1 cara.)}

\begin{exercisewrap}
\begin{nexercise}
Ada berapa cara untuk menempatkan satu sukses dan $n-1$ kegagalan dalam
$n$ percobaan? Ada berapa cara untuk menempatkan $n-1$ sukses dan satu
kegagalan dalam $n$ percobaan?\footnotemark{}
\end{nexercise}
\end{exercisewrap}
\footnotetext{Untuk satu sukses dan $n-1$ kegagalan, terdapat tepat $n$
  posisi berbeda tempat kita dapat menaruh sukses. Jadi, terdapat $n$
  cara untuk menempatkan satu sukses dan $n-1$ kegagalan. Argumen serupa
  digunakan untuk pertanyaan kedua. Secara matematis, kita menunjukkan
  hasil-hasil ini dengan memverifikasi kedua persamaan berikut:
  \begin{align*}
  {n \choose 1} = n,
    \qquad {n \choose n-1} = n
  \end{align*}}


\newpage


\subsection{Pendekatan normal terhadap distribusi binomial}
\label{normalApproxBinomialDistSubsection}

\index{distribution!binomial!normal approximation|(}

Rumus binomial merepotkan ketika ukuran sampel ($n$) besar, terutama
ketika kita meninjau suatu rentang pengamatan. Dalam beberapa kasus,
kita dapat menggunakan distribusi normal sebagai cara yang lebih mudah
dan cepat untuk menaksir peluang binomial.

\newcommand{\smokeprop}{0.15}
\newcommand{\smokeperc}{15\%}
\newcommand{\smokepropcomp}{0.85}
\newcommand{\smokeperccomp}{85\%}
\newcommand{\smokex}{42}
\newcommand{\smokexplusone}{43}
\newcommand{\smoken}{400}
\newcommand{\smokelowertailbinom}{0.0054}
\newcommand{\smokemean}{60}
\newcommand{\smokemeancomp}{340}
\newcommand{\smokesd}{7.14}
\newcommand{\smokez}{-2.52}
\newcommand{\smokelowertailnormal}{0.0059}

\begin{examplewrap}
\begin{nexample}{Sekitar \smokeperc{} penduduk Amerika Serikat merokok.
    Pemerintah daerah meyakini bahwa proporsi perokok di komunitasnya
    lebih rendah, lalu meminta pelaksanaan survei terhadap 400 orang yang
    dipilih secara acak. Survei tersebut menemukan bahwa hanya \smokex{} dari
    \smoken{} peserta yang merokok. Jika proporsi perokok yang sebenarnya
    di komunitas itu memang \smokeperc{}, berapa peluang untuk mengamati
    \smokex{} perokok atau kurang dalam sampel berisi \smoken{} orang?}
  \label{exactBinomSmokerExSetup}%
  Verifikasi seperti biasa bahwa keempat syarat model binomial berlaku
  kita serahkan sebagai latihan.

  Pertanyaan yang diajukan setara dengan menanyakan peluang untuk
  mengamati $k=0$, 1, 2, ..., atau \smokex{} perokok dalam sampel berisi
  $n=\smoken{}$ orang ketika $p=\smokeprop{}$. Kita dapat menghitung
  \smokexplusone{} peluang yang berbeda ini dan menjumlahkannya untuk
  memperoleh jawaban:
  \begin{align*}
  &P(k=0\text{ atau }k=1\text{ atau }\cdots\text{ atau } k=\smokex{}) \\
	&\qquad = P(k=0) + P(k=1) + \cdots + P(k=\smokex{}) \\
	&\qquad = \smokelowertailbinom{}
  \end{align*}
  Jika proporsi perokok yang sebenarnya di komunitas adalah
  $p=\smokeprop{}$, peluang untuk mengamati \smokex{} perokok atau kurang
  dalam sampel berukuran $n=\smoken{}$ adalah \smokelowertailbinom{}.
\end{nexample}
\end{examplewrap}

Perhitungan dalam Contoh~\ref{exactBinomSmokerExSetup} membosankan dan
panjang. Secara umum, kita sebaiknya menghindari pekerjaan semacam itu
jika tersedia metode alternatif yang lebih cepat, lebih mudah, dan tetap
akurat. Ingat bahwa menghitung peluang suatu rentang nilai jauh lebih
mudah dengan model normal. Kita mungkin bertanya-tanya apakah model normal
layak digunakan sebagai pengganti distribusi binomial. Ternyata jawabannya
ya, asalkan syarat-syarat tertentu terpenuhi.

\begin{exercisewrap}
\begin{nexercise}
Di sini kita meninjau model binomial ketika peluang sukses adalah
$p=0.10$. Gambar~\ref{fourBinomialModelsShowingApproxToNormal}
menampilkan empat histogram berongga untuk sampel simulasi dari distribusi
binomial dengan empat ukuran sampel berbeda: $n=10, 30, 100, 300$.
Apa yang terjadi pada bentuk distribusi ketika ukuran sampel meningkat?
Distribusi apakah yang menyerupai histogram berongga terakhir?\footnotemark{}
\end{nexercise}
\end{exercisewrap}
\footnotetext{Distribusinya berubah dari distribusi berbentuk balok dan
  menceng menjadi distribusi yang cukup menyerupai distribusi normal pada
  histogram berongga terakhir.}

\begin{figure}[h]
  \centering
  \Figure[Ditampilkan empat histogram berongga, masing-masing pada grafik tersendiri, berdasarkan peluang p sama dengan 0.10 dan ukuran sampel n sama dengan 10, 30, 100, dan 300. Grafik pertama untuk n = 10 menunjukkan distribusi yang berpusat di 1 dan sangat menceng kanan. Grafik kedua untuk n = 30 menunjukkan distribusi yang berpusat di sekitar 3, sedikit menceng kanan, dan mulai tampak agak menyerupai distribusi berbentuk lonceng. Grafik ketiga untuk n = 100 menunjukkan distribusi yang berpusat di sekitar 10 dan nyaris sepenuhnya simetris, dengan hanya sedikit tanda kemencengan kanan. Distribusi ketiga ini juga tampak sangat menyerupai lonceng. Grafik keempat untuk n = 300 menunjukkan distribusi yang berpusat di sekitar 30 dan simetris. Grafik terakhir ini tampak sangat menyerupai lonceng dan menyerupai distribusi normal.]{0.92}{fourBinomialModelsShowingApproxToNormal}
  \caption{Histogram berongga sampel dari model binomial ketika
      $p=0.10$. Ukuran sampel keempat grafik tersebut berturut-turut
      adalah $n=10$, 30, 100, dan 300.}
  \label{fourBinomialModelsShowingApproxToNormal}
\end{figure}

\begin{onebox}{Pendekatan normal terhadap distribusi binomial}
  Distribusi binomial dengan peluang sukses $p$ mendekati distribusi
  normal apabila ukuran sampel $n$ cukup besar sehingga $np$ dan
  $n(1-p)$ masing-masing sekurang-kurangnya 10. Distribusi normal yang
  digunakan sebagai pendekatan memiliki parameter yang bersesuaian dengan
  rata-rata dan simpangan baku distribusi binomial:\vspace{-1.5mm}
  \begin{align*}
  \mu &= np
      &\sigma& = \sqrt{np(1 - p)}
  \end{align*}
\end{onebox}

Pendekatan normal dapat digunakan ketika menghitung peluang bagi rentang
yang mencakup banyak kemungkinan banyaknya sukses. Sebagai contoh, kita
dapat menerapkan
distribusi normal pada latar Contoh~\ref{exactBinomSmokerExSetup}.

\D{\newpage}

\begin{examplewrap}
\begin{nexample}{Bagaimana kita dapat menggunakan pendekatan normal untuk
    menaksir peluang mengamati \smokex{} perokok atau kurang dalam sampel
    berukuran \smoken{}, jika proporsi perokok yang sebenarnya adalah
    $p=\smokeprop{}$?}
  \label{approxNormalForSmokerBinomEx}
  Menunjukkan bahwa model binomial layak digunakan merupakan latihan yang
  disarankan dalam Contoh~\ref{exactBinomSmokerExSetup}. Kita juga
  memverifikasi bahwa $np$ dan $n(1-p)$ masing-masing sekurang-kurangnya 10:
  \begin{align*}
  np &= \smoken{} \times \smokeprop{} = \smokemean{}
  &n (1 - p) &= \smoken{} \times \smokepropcomp{}
      = \smokemeancomp{}
  \end{align*}
  Setelah syarat-syarat ini diperiksa, kita dapat menggunakan pendekatan
  normal sebagai pengganti distribusi binomial dengan rata-rata dan
  simpangan baku dari model binomial:
  \begin{align*}
  \mu &= np = \smokemean{}
  &\sigma &= \sqrt{np(1 - p)} = \smokesd{}
  \end{align*}
  Kita ingin mencari peluang mengamati \smokex{} perokok atau kurang
  dengan model ini.
\end{nexample}
\end{examplewrap}

\begin{exercisewrap}
\begin{nexercise}
Gunakan model normal $N(\mu=\smokemean{}, \sigma=\smokesd{})$ untuk
menaksir peluang mengamati \smokex{} perokok atau kurang. Jawaban Anda
seharusnya kira-kira sama dengan penyelesaian
Contoh~\ref{exactBinomSmokerExSetup}:%
~\smokelowertailbinom{}.~\footnotemark{}
\end{nexercise}
\end{exercisewrap}
\footnotetext{Hitung skor-Z terlebih dahulu:
  $Z = \frac{\smokex{} - \smokemean{}}{\smokesd{}} = \smokez{}$.
  Luas ekor kiri yang bersesuaian adalah \smokelowertailnormal{}.}



\newpage


\subsection{Pendekatan normal tidak akurat pada interval kecil}

Pendekatan normal terhadap distribusi binomial cenderung memberikan
taksiran yang buruk ketika menaksir peluang bagi rentang kecil banyaknya
sukses, bahkan apabila syarat-syaratnya terpenuhi.

\newcommand{\smokeA}{49}
\newcommand{\smokeB}{50}
\newcommand{\smokeC}{51}
\newcommand{\smokeABCBinom}{0.0649}
\newcommand{\smokeABCNormal}{0.0421}
\newcommand{\smokeABCNormalFixed}{0.0633}

Misalkan kita ingin menghitung peluang mengamati \smokeA{}, \smokeB{},
atau \smokeC{} perokok di antara \smoken{} orang ketika
$p=\smokeprop{}$. Dengan sampel sebesar itu, kita mungkin tergoda untuk
menerapkan pendekatan normal dan menggunakan rentang \smokeA{} hingga
\smokeC{}. Namun, kita akan menemukan bahwa penyelesaian binomial dan
pendekatan normal berbeda secara nyata:
\begin{align*}
\text{Binomial:}&\ \smokeABCBinom{}
&\text{Normal:}&\ \smokeABCNormal{}
\end{align*}
Kita dapat mengidentifikasi penyebab perbedaan ini dengan menggunakan
Gambar~\ref{normApproxToBinomFail}, yang menampilkan daerah yang mewakili
peluang binomial (dibatasi garis) dan pendekatan normal (diarsir).
Perhatikan bahwa daerah di bawah distribusi normal lebih sempit sebesar
0.5 satuan pada kedua sisi interval.

\begin{figure}[h]
  \centering
  \Figure[Ditampilkan sebuah distribusi normal yang berpusat di 60 dengan simpangan baku sekitar 7. (Penentuan bahwa simpangan bakunya sekitar 7 didasarkan pada distribusi normal yang sangat mendekati 0 pada jarak sekitar 20 dari rata-rata; hal ini terjadi pada jarak sekitar 3 simpangan baku dari rata-rata.) Suatu daerah pada distribusi ini diarsir dari 49 hingga 51. Selain itu, sebuah daerah dengan garis tepi merah membentang dari 48.5 hingga 51.5 untuk mewakili distribusi binomial eksak.]{1.0}{normApproxToBinomFail}
  \caption{Kurva normal dengan daerah antara \smokeA{} dan \smokeC{}
      diarsir. Daerah yang dibatasi garis mewakili peluang binomial eksak.}
  \label{normApproxToBinomFail}
\end{figure}

\begin{onebox}{Memperbaiki pendekatan normal terhadap distribusi binomial}
  Pendekatan normal terhadap distribusi binomial pada interval nilai biasanya
  menjadi lebih baik jika nilai batasnya sedikit diubah. Nilai batas pada
  ujung bawah daerah yang diarsir perlu dikurangi 0.5, sedangkan nilai
  batas pada ujung atas perlu ditambah 0.5.
\end{onebox}

Anjuran untuk menambahkan luas tambahan ketika menerapkan pendekatan
normal paling sering berguna saat menelaah suatu rentang pengamatan.
Dalam contoh di atas, taksiran distribusi normal yang telah diperbaiki
adalah \smokeABCNormalFixed{}, jauh lebih dekat dengan nilai eksak
\smokeABCBinom{}. Meskipun koreksi ini juga dapat diterapkan ketika
menghitung luas ekor, manfaat koreksi tersebut biasanya hilang karena
keseluruhan intervalnya umumnya cukup lebar.

\index{distribution!binomial!normal approximation|)}
\index{distribution!binomial|)}


{\input{ch_distributions/TeX/binomial_distribution.tex}}




%_________________
